\begin{center}\begin{tabular}{rrr}\toprule
步长 / s & 单摆周期 / s & 10 s 最大相对能量漂移 \\\midrule
0.0200 & 2.03786885 & $3.891\times10^{-7}$ \\
0.0100 & 2.03786764 & $1.240\times10^{-8}$ \\
0.0050 & 2.03786785 & $4.034\times10^{-10}$ \\
0.0025 & 2.03786791 & $1.360\times10^{-11}$ \\
\bottomrule\end{tabular}\end{center}
独立 Simpson 积分有限振幅周期为 2.03786792 s, 小角周期为 2.00606668 s. 周期插值本身只有二阶, 不能把过零时间误差也称为四阶.
链条采用前述无耗散折返模型, $L=1$ m. 用 $y=L\sin^2\phi$ 消去端点根号, 实际积分得如下支力曲线点; 不把末端奇异力当成实测峰值.
\begin{center}\begin{tabular}{rrr}\toprule
$y/L$ & $t$ / s & $F_S/(Mg)$ \\\midrule
0.00 & 0.000000 & 0.0000 \\
0.10 & 0.200179 & 0.2056 \\
0.25 & 0.311917 & 0.5417 \\
0.50 & 0.428056 & 1.2500 \\
0.75 & 0.502842 & 2.6250 \\
0.90 & 0.531449 & 5.8500 \\
0.99 & 0.540688 & 50.9850 \\
\bottomrule\end{tabular}\end{center}
折返完成时刻的有限极限为 0.54098864 s; 此时理想支力发散, 最终静止的 $Mg$ 不是该曲线连续延拓.
另运行期末第 7 题的中心力初值: $A=m=L=1$, 初 $r=1.01$, 纯切向速率 $1/1.01$. 步长 0.001, 积分 20 时间单位. 实得径向周期 2.56495053, 线性预测 $2\pi/\sqrt6=2.56509966$; 有限振幅差约 5.81e-05. 最大相对能量漂移 1.03e-14. 数据仅验证这一接近圆轨道的复现, 一般结论由正文推导保证.
